\section{The constant-plus-power boundary}
\label{app:constant-proof}

Write \(F_{n,d}(t)=\sum_{k=0}^n\binom nk(dk)!t^k\) and
\(R_n(z)=z^nF_{n,d}(1/z)\). We prove that \(R_n\) is squarefree
for all positive integers \(n,d\). The proof uses a coefficient
bound for large \(d\), Newton residuals for the small prime powers,
and a uniform contour estimate for the remaining large orders.

\subsection{The coefficient bound}

For integers \(1\le u\le v\) one has
\[
 \frac{u!}{(v!)^{u/v}}\le (u/v)^{u/2}.
\]
Indeed \((k!)^{1/k}/\sqrt{k}\) is nondecreasing. The comparison
at \(k,k+1\), after raising to positive powers, is
\((k!)^2(k+1)^{k(k-1)}\le k^{k(k+1)}\). This follows from
\(k!\le((k+1)/2)^k\) and \((1+1/k)^{k+1}\le4\). For the latter
bound, \((x+1)\log(1+1/x)\) decreases for \(x>0\).
Also \(k\log(n/k)\ge(\log2)\min(k,n-k)\), by concavity of
\(-x\log x\) on \([1/2,1]\) and by \(n/k\ge2\) for \(k\le n/2\).

After scaling \(t\) by \(((dn)!)^{-1/n}\), the polynomial is
\(1+t^n+\sum_{k=1}^{n-1}b_kt^k\), with
\[
 0<b_k\le C_k:=\binom nk q^{\min(k,n-k)},\qquad q=2^{-d/2}.
\]
The bound follows from the factorial inequality at \(dk,dn\).
If \(2^d\ge4n^2\), then \(q\le1/(2n)\).
By symmetry of \(C_k\) and \(q^{\min(k,n-k)}\le q^k+q^{n-k}\),
both weighted sums \(\sum(k/n)b_k\) and
\(\sum(1-k/n)b_k\) are bounded by \((1+q)^n-1<1\).
At a repeated root of modulus at most one, subtract \(t/n\) times
the derivative from the polynomial to obtain a contradiction to
the second bound. At a repeated root of modulus at least one,
divide \(t/n\) times the derivative by \(t^n\) and use the first.
We have proved that
\begin{equation}
 2^d\ge4n^2 \ \Longrightarrow\  F_{n,d}\text{ is squarefree}.
 \label{eq:constant-large-gap}
\end{equation}

\subsection{Small exponents}

For \(d=1\), \(R_n(z)=n!\sum_{j=0}^n z^j/j!\). Its derivative
and itself cannot share a root, since the two truncated
exponentials differ by \(z^n/n!\). To cover \(2\le d\le5\),
we give the prime-power argument. The factorial units and the
digit structure of the Newton polygon used here occur in
Rocks's work on factorial moment polynomials~\cite{rocks}.

Let \(d=p^\nu\) with \(p\) prime and \(\nu\ge1\),
and put \(c=(d-1)/(p-1)\). For an integer \(u\ge0\), write
\(U_p(u!)=p^{-v_p(u!)}u!\). If \(u=\sum_i u_i p^i\),
Wilson's theorem, applied after taking out multiples of \(p\)
one level at a time, gives
\[
 U_p(u!)\equiv(-1)^{v_p(u!)}\prod_i u_i!\pmod p.
\]
The digit shift from \(k\) to \(dk\) and Legendre's formula
therefore give
\[
 F_{n,d}(T/p^c)=\sum_{k=0}^n \frac{n!}{(n-k)!}w_k T^k,\qquad
 w_k=\frac{(dk)!}{p^{ck}k!}\in\Z_p^\times,\qquad
 w_k\equiv(-1)^{ck}\pmod p .
\]
Write \(n=\sum_i a_i p^i\) and \(s_i=\sum_{j\le i}a_jp^j\),
with \(s_{-1}=0\). The lower Newton polygon has an edge from
\(s_{i-1}\) to \(s_i\) for each \(a_i>0\), of slope \(h_i/e_i\)
with \(e_i=p^i\) and \(h_i=(p^i-1)/(p-1)\). Indeed the
coefficient ordinate is
\(v_p(n!)-v_p((n-k)!)\). For \(k=s_i-\delta\) in that segment
the difference between the ordinate and the line through
the two endpoints is
\[
 \frac{\operatorname{digitsum}_p(\delta)-\delta/p^i}{p-1}\ge0.
\]
This is Legendre's formula, since \(n-s_i\) and \(\delta\)
have disjoint digit positions. Equality occurs exactly at the
indices spaced by \(e_i\), the slopes strictly increase,
and \(h_i/e_i\) is in lowest terms.

Divide each coefficient on the edge by its power of \(p\)
and take ratios to the leftmost coefficient. Wilson's unit
formula then gives the edge polynomial
\[
 \sum_{j=0}^{a_i}\frac{a_i!}{(a_i-j)!}(\eta_i Y)^j,\qquad
 \eta_i=(-1)^{h_i+c e_i}.
\]
Here only digit \(a_i\) changes in the factorial ratio, and the
change in valuation is \(jh_i\). The displayed polynomial
is a scaled reversal of \(E_{a_i}(Z)=\sum_{j=0}^{a_i}Z^j/j!\)
over \(\mathbb F_p\). It is squarefree because
\(E_{a_i}'=E_{a_i-1}\) and \(E_{a_i}-E_{a_i-1}=Z^{a_i}/a_i!\).

For completeness, edge polynomials multiply at any fixed
slope \(u/e\) in lowest terms. Substitute \(T=\pi^{-u}Z\)
in a valued extension with \(\pi^e=p\), divide each polynomial
by the power giving minimum coefficient valuation zero,
and reduce. The reductions multiply. The surviving indices
in a polynomial over \(\Q_p\) are all congruent modulo \(e\),
so after removing the lowest monomial the same identity
holds in the edge variable \(Y=Z^e\).
A square divisor would thus give a repeated factor in the
edge polynomial at any one of its slopes. The squarefreeness
above rules this out.

\subsection{A uniform contour estimate}

We now assume \(n\ge10000\) and \(6\le d\le\sqrt n\).
Set \(A=(d/e)^d\) and \(v=z/(n^d A)\).
The exponential integral for \(R_n\), with its contour split
at \(d/2\), is
\[
 R_n(n^d A v)=n^{dn+1}I(v),\qquad
 I(v)=\int_0^\infty(Av+s^d)^n e^{-ns}\,ds=K(v)+J(v),
\]
where \(K\) and \(J\) integrate over \([0,d/2]\) and
\([d/2,\infty)\), respectively. Both are polynomials in \(v\).
We bound \(J\) by a Gaussian at its saddle and \(K\) by an
endpoint approximation. Their logarithmic derivatives are
separated where a zero of \(I\) can occur.

\paragraph{The outer bound.}
For \(\theta=k/n\), the entropy bound on \(\binom nk\) and
\((dk)!\le e\,dk(dk/e)^{dk}\) give
\[
 \frac{[\binom nk(dk)!]^{1/k}}{n^d A}
 \le q_d(\theta)(e\,dk)^{1/k},\qquad
 q_d(\theta)=\theta^{d-1}(1-\theta)^{-(1-\theta)/\theta}.
\]
The endpoint values use continuous limits. We have \(q_d\le q_6\).
The sign of the derivative of \(\log q_6\) is the sign of
\(6\theta+\log(1-\theta)\). Its positive-to-negative change
occurs at \(\theta=1-u\) with \(0<u<1/300\), since
\(6(299/300)>\log300\).
There \(\log q_6=5\log(1-u)+6u<u\), so \(q_6<300/299\).
If \(k\ge n/10\), then \(d\le\sqrt n\) gives
\[
 k^{-1}\log(e\,dk)\le10(1+\tfrac32\log n)/n<2/125.
\]
Thus the normalized coefficient root is less than
\((300/299)(125/123)<21/20\).
If \(k<n/10\), the bounds
\(\binom nk^{1/k}\le en/k\) and \((e\,dk)^{1/k}\le ed\)
give instead \(e^2 d\theta^{d-1}<21/20\).
The latter unit bound uses \(k\le2^{k-1}\le(ed)^{k-1}\).
Summing the resulting finite geometric series excludes
\(|v|\ge R=21/10\).

\paragraph{The large saddle.}
Put \(B=11/5\) and \(\delta=1/25\). For \(|v|\le B\) the
equation \(A v=s^{d-1}(d-s)\) has a unique solution
\(s_0=\alpha_0+i\beta_0\) satisfying \(|s_0-d|<\delta\).
On \(|s-d|=\delta\) the modulus on the right divided by \(A\)
is at least \(\delta(1-\delta)e^d/d>B\).
Bernoulli's inequality gives the lower estimate, and \(e>27/10\)
verifies the strict comparison at \(d=6\), where the bound
is smallest. Rouch\'e's theorem gives the solution, which
is simple and analytic in \(v\) on a neighborhood of the disk.
Use the phase
\[
 \Phi(s,v)=d\log s+\log(1+A v/s^d)-s,\qquad
 S=\Re\Phi(s_0,v),\qquad a_0=1-(d-1)/s_0 .
\]
At \(s_0\), \(\Phi_s=0\), \(\Phi_{ss}=-a_0\), and
\begin{equation}
 \Re a_0>9/(10d),\qquad |a_0|<11/(10d),\qquad
 S>\log A-1/100 .
 \label{eq:constant-saddle-bounds}
\end{equation}
For the first two bounds use
\(|a_0-1/d|=(d-1)|s_0-d|/(d|s_0|)<\delta/d\).
For the height, use
\(\Phi(s_0,v)=(d-1)\log s_0+\log d-s_0\).
The loss is at most \(\delta/d+\delta^2/[2d(1-\delta/6)]<1/100\)
by the logarithm series.

Deform the contour of \(J\) to its horizontal translate
by \(i\beta_0\), with the vertical connector at \(d/2\).
The integrand is entire and the connector at positive
infinity tends to zero. On \(s=s_0+t,\ |t|\le d/8\)
we have \(|s|>6d/7\). With \(q_t=B(35/81)^6\),
differentiating the convergent logarithm series gives
\[
 d^2|\Phi_{sss}|
 \le (7/6)^3\left[2+
 \frac{36q_t(1+q_t)}{(1-q_t)^3}
 +\frac{18q_t}{(1-q_t)^2}+\frac{2q_t}{1-q_t}\right]<8.
\]
For this inequality use
\(|Av/s^d|\le B(35/81)^d\).
In the differentiated series the corresponding terms have
factors \(d^2,d,1\), each giving its largest bound at \(d=6\).
The cubic remainder is at most \(4|t|^3/(3d^2)\).
By \eqref{eq:constant-saddle-bounds} its contribution to the local integral is at most
\[
 \frac{4ne^{nS}}{3d^2}\int_{\R}|t|^3 e^{-17nt^2/(60d)}\,dt
 <17e^{nS}/n.
\]
Completing the Gaussian costs less than
\(18e^{nS-9dn/1280}/n\), by bounding the Gaussian tail by
the integral with the extra factor \(8|t|/d\).
The complete value
\(U_0=e^{n\Phi(s_0,v)}\sqrt{2\pi/(na_0)}\)
has modulus greater than \(2\sqrt{d/n}\,e^{nS}\).

We give the bounds on the remaining pieces of this contour.
For \(2d/3\le x\le4d/3\), let \(q_c=B(5/9)^6\).
On \(s=x+i\beta_0\), the logarithm series and
\(|\beta_0|/x<1/100\) give
\[
 \frac{d^2}{dx^2}\Re\Phi(x+i\beta_0,v)
 \le-\frac d{x^2}\left[
 \frac{1-10^{-4}}{(1+10^{-4})^2}
 -\frac{6q_c}{(1-q_c)^2}-\frac{q_c}{1-q_c}\right]
 <-\frac3{16d}.
\]
Here \(|Av/s^d|\le B(5/9)^d\), and \(d(5/9)^d\) decreases
for \(d\ge6\). Integrating from \(x=\alpha_0\) bounds the
height by \(S-3(x-\alpha_0)^2/(32d)\).
Outside the local interval this costs at least \(d/1000\)
in height, over length less than \(d\).
At \(4d/3\) the loss is greater than \(d/100\).
For \(x\ge4d/3\), the real derivative is at most
\((3/4)/(1-B(5/18)^6)-1<-1/5\), so the far ray costs at
most \(5e^{nS-dn/100}/n\).

On \([d/2,2d/3]\) the upper envelope
\([BA+(x+\delta)^d]e^{-x}\) increases. For
\(x+\delta<d-1\), the function
\((x+\delta)^{d-1}(d-x-\delta)\) increases, and its value
at \(x=d/2\) exceeds \(BA\) since \((e/2)^d>B\).
The envelope divided by \(A\) is at most
\(B e^{-2d/3}+r_\ell^d\) there, with
\(r_\ell=(101/150)(141/100)\) from \(e^{1/3}<141/100\).
The corresponding bound on the vertical connector is
\(B e^{-d/2}+r_v^d\) with \(r_v=38/45\), from \(e^{1/2}<5/3\).
Both bounds are below \(e^{-d/50}\).
After multiplication by \(e^{d/50}\) their two exponential
terms decrease in \(d\), and at \(6\) the comparisons follow
from
\[
 B(10/27)^3+r_\ell^6/(22/25)<1,\qquad
 B(10/27)^2+r_v^6/(22/25)<1.
\]
By \eqref{eq:constant-saddle-bounds} the height of these pieces is below \(S-d/100\).
Their lengths are \(d/6\) and at most \(\delta\).
The phase charts used on the horizontal contour are valid
since \(|Av/s^d|\le B(2/e)^d<1\) for \(x\ge d/2\).

Summing the contour bounds gives an error at most \(e^{nS}\)
times
\[
 40/n+d e^{-dn/1000}+(d+\delta)e^{-dn/100}<44/n .
\]
The last comparison uses \(e^x\ge x^3/6\) in the middle
term and \(e^x\ge x^2/2\) in the last.
Consequently \(J=U_0(1+\epsilon)\) with
\(|\epsilon|<22/\sqrt{dn}<1/10\).
The error is analytic on a neighborhood of the disk.
Cauchy's estimate at radius \(1/10\) gives
\(|\epsilon'|<1\) on \(|v|\le R\). Also
\[
 U_0'/U_0=n\frac{d-s_0}{dv}-\frac{a_0'}{2a_0},\qquad
 s_0'=\frac{A}{d s_0^{d-2}(d-1-s_0)} .
\]
Bernoulli's inequality gives \(|s_0'|\le de^{-d}/(1-\delta)^2<1\),
and hence \(|a_0'/(2a_0)|<1\) by \eqref{eq:constant-saddle-bounds}. We obtain
\begin{equation}
 \left|J'/J-n(d-s_0)/(dv)\right|<3 \qquad (|v|\le R).
 \label{eq:constant-large-derivative}
\end{equation}
The expression at zero is understood by continuity.

\paragraph{The endpoint.}
For \(0\le x\le d/2\),
\[
 \log(1+x^d/((2/5)A))\le19x/20.
\]
At \(x=d/2\) this is \(1+(5/2)(e/2)^d<e^{19d/40}\).
After multiplication by \(e^{-19d/40}\), both terms
decrease with \(d\), since \(e^{21/40}\le(5/3)(40/39)<2\).
At \(d=6\) use \(e<49/18\) and the rational comparison
\[
 1+(5/2)(49/36)^6<\sum_{j=0}^{12}(57/20)^j/j! .
\]
To pass to smaller \(x\), the logarithm divided by \(x\)
has derivative of the sign of \(du/(1+u)-\log(1+u)\),
where \(u=x^d/((2/5)A)\). This function is nonnegative up
to \(u=d-1\) by differentiation. Thereafter its first term
is at least \(d-1>19d/40\), while the logarithm is bounded
by its endpoint value.

For \(|v|\ge2/5\), write \(K=(Av)^n(1+D(v))/n\).
Subtract the full exponential integral and use
\(|(1+z)^n-1|\le n|z|(1+|z|)^{n-1}\).
For \(D'\) differentiate \((1+x^d/(Av))^n\) inside
the integral. The same endpoint envelope gives
\[
 |D(v)|\le \frac{d!20^{d+1}}{(2/5)A n^{d-1}}+e^{-dn/2},\qquad
 |D'(v)|\le \frac{d!20^{d+1}}{(2/5)^2 A n^{d-1}}.
\]
Both are below \(1/1000\), since \(d!\le ed A\), with the
largest majorants at \(d=6,n=10000\). In particular
\begin{equation}
 K(v)\ne0,\qquad |K'/K-n/v|<1\qquad(|v|\ge2/5).
 \label{eq:constant-endpoint-derivative}
\end{equation}

For \(|v|\le1/2\), the upper envelope
\((A/2+x^d)e^{-x}\) attains its maximum on \([0,d/2]\)
at an endpoint, by the sign of its derivative.
At the right endpoint it is below \(A/2\) because
\(e^{-d/2}+2(e^{1/2}/2)^d\le(5/8)^6+2(5/6)^6<1\).
Thus \(|K|\le(d/2)(A/2)^n\).
The lower estimate for \(J\) gives
\(|J|>(3/2)\sqrt{d/n}(A e^{-1/100})^n\).
It follows that \(|K/J|<\sqrt{dn}(50/99)^n/3<1\),
using \(\sqrt{dn}\le n\) and \(\sum_{j\ge1}j(50/99)^j<3\).
Therefore every zero of \(I\) lies in \(1/2<|v|<R\).
At such a zero, \(I'=J(J'/J-K'/K)\), and
\eqref{eq:constant-large-derivative} and
\eqref{eq:constant-endpoint-derivative} give
\[
 |J'/J-K'/K|>n|s_0|/(d|v|)-4>n/3-4>0 .
\]
This proves simplicity in the entire stated contour range.

\subsection{Completion and finite verification}

If \(n\ge10000\) and \eqref{eq:constant-large-gap} fails,
then \(d<2\log_2(2n)<\sqrt n\). The second inequality
follows since \(u-2-4\log_2 u\) is positive at \(u=100\)
and increasing for \(u\ge100\). Hence the contour estimate
and the small-exponent argument cover every such order.
For \(n<10000\), the remaining pairs have \(6\le d\le29\)
and \(2^d<4n^2\). The case \(n=1\) is linear.

Here is a small-degree test for that finite range.
Let \(n=p+r\) with \(p\) prime and \(p>d\max(1,r)\).
Then \(r<p\) and \(p^2-dn=p(p-d)-dr>0\).
The Newton ordinate of \(F_{n,d}\) is
\[
 y(k)=\lfloor dk/p\rfloor+\mathbf1_{\{k\bmod p>r\}}.
\]
It is zero up to \(r\). For \(r<k<p\) it exceeds \(dk/p\),
and for \(p\le k\le n\) it equals \(d\).
Thus the polygon has a horizontal part of width \(r\)
and a slope \(d/p\) of width \(p\).
Every repeated rational irreducible divisor has degree
at most \(n/2<p\), so by the slope denominator it can
use only the horizontal part. It is monic integral in
the reciprocal and reduces with nonzero constant term.
Lucas's congruence gives \(R_n\bmod p=z^p(R_r\bmod p)\).
Indeed for \(k\le r\) the binomial coefficient reduces
to \(\binom rk\), and the other coefficients of \(F_{n,d}\)
vanish modulo \(p\). Thus a repeated divisor would
reduce to a repeated divisor of \(R_r\bmod p\).

Every repeated rational irreducible divisor has degree
at least two. For even \(n\), the exponential integral
\(F_{n,d}(t)=\int_0^\infty(1+tu^d)^ne^{-u}\,du\) is positive
on the real axis. For odd \(n\), its derivative is positive there.
In either case it has no repeated real root.
It follows that a gcd of degree at most one between the modular
\(R_r\) and its derivative excludes every repeated rational
factor of \(R_n\). The same argument applies to a direct
modular test of \(R_n\) at any prime greater than \(dn\).

The supplementary verifier exhausts \(202\,021\) pairs in
the remaining range. It uses a shifted prime in \(201\,927\)
cases, where the tested degree is at most \(35\). In the
other \(94\) cases it tests the full polynomial at
\(1\,000\,000\,007\), with degree at most \(148\). The primes
are obtained by an Eratosthenes sieve in the shifted tests
and by trial division for the full tests. Every modular gcd
has degree at most one. This completes the squarefreeness
proof and thus Lemma~\ref{lem:constant-support}.
